% Part: computability
% Chapter: recursive-functions
% Section: bounded-minimization

\documentclass[../../../include/open-logic-section]{subfiles}

\begin{document}

\olfileid{cmp}{rec}{bmi}
\olsection{સીમાબદ્ધ લઘુત્તમીકરણ}

\begin{explain}
કોઈ ગુણધર્મ અથવા સંબંધ~$P$ સંતોષતી સૌથી નાની સંખ્યા તરીકે વિધેય
વ્યાખ્યાયિત કરવું ઘણી વાર ઉપયોગી છે. $P$નો નિર્ણય કરી શકાતો હોય,
તો $0$, $1$, $2$, \dots{} બધી શક્ય સંખ્યાઓ ક્રમે અજમાવતાં~$P$
સંતોષતી સૌથી નાની સંખ્યા મળે અને આ વિધેય ગણી શકાય. આવી અનિયત
મર્યાદા સુધીની શોધ આપણને આદિમ પુનરાવર્તી વિધેયોના ક્ષેત્રની બહાર લઈ
જાય છે. પરંતુ જો આપણે \emph{સ્વતંત્ર રીતે આપેલી કોઈ સીમાથી નાની}
સૌથી નાની સંખ્યા જ શોધવી હોય, તો આદિમ પુનરાવર્તી ક્ષેત્રમાં રહીએ
છીએ. બીજા શબ્દોમાં, અને થોડું વધુ સામાન્ય રીતે, ધારો કે
સંબંધ~$R(x,z)$ આદિમ પુનરાવર્તી છે. $x$ અને~$y$ને એવા સૌથી નાના
$z < y$ પર મોકલતું વિધેય વિચારો કે જેના માટે $R(x, z)$ લાગુ પડે.
$R(x, 0)$, $R(x, 1)$, \dots, $R(x, y-1)$ની તપાસ કરીને તેને પણ
ગણી શકાય. પણ તે આદિમ પુનરાવર્તી શા માટે છે?
\end{explain}

\begin{prop}
જો $R(\vec x, z)$ આદિમ પુનરાવર્તી હોય, તો વિધેય
$m_R(\vec{x}, y)$ પણ આદિમ પુનરાવર્તી છે. $R(\vec x, z)$ લાગુ પડે
એવી કોઈ સંખ્યા $y$થી નાની હોય, તો આ વિધેય સૌથી નાનો એવો~$z$ આપે
છે; અન્યથા $y$ આપે છે. વિધેય~$m_R$ને આ રીતે લખીશું:
\[
\bmin{z < y}{R(\vec{x}, z)},
\]
\end{prop}

\begin{proof}
$z < 0$ એવી સંખ્યા જ નથી; તેથી $R(\vec{x}, z)$ લાગુ પડે
એવી~$z < 0$ પણ હોઈ શકે નહીં. એટલે $m_R(\vec x, 0) = 0$.

સીમા $y + 1$ સ્વરૂપની હોય ત્યારે ત્રણ કિસ્સા છે:
\begin{enumerate}
\item $R(\vec{x}, z)$ લાગુ પડે એવી $z < y$ છે. ત્યારે
$m_R(\vec{x}, y+1) = m_R(\vec{x}, y)$.
\item એવી કોઈ~$z<y$ નથી, પરંતુ $R(\vec{x}, y)$ લાગુ પડે છે.
ત્યારે $m_R(\vec{x}, y+1) = y$.
\item $R(\vec{x}, z)$ લાગુ પડે એવી $z < y+1$ નથી. ત્યારે
$m_R(\vec{x}, y+1) = y+1$.
\end{enumerate}
\sourcecorrection{OLCMP-004}{સ્થિર અંગ્રેજી મૂળના ત્રીજા કિસ્સામાં
ડાબી બાજુ $m_R(\vec z, y+1)$ છપાયું છે. બાકીના બે કિસ્સા, પ્રસ્તાવનું
વિધેય અને તરત નીચેનું સમીકરણ બધે $\vec x$ છે; ગુજરાતી ત્રીજા
કિસ્સામાં પણ $m_R(\vec{x}, y+1)$ રાખ્યું છે.}
આથી $m_R(\vec x, 0)$ને આદિમ પુનરાવર્તન વડે નીચે મુજબ વ્યાખ્યાયિત
કરી શકીએ:
\begin{align*}
m_R(\vec{x}, 0) & = 0\\
m_R(\vec{x}, y+1) & =
\begin{cases}
m_R(\vec{x}, y) & \text{જો $m_R(\vec x, y) \neq y$ હોય}\\
y & \text{જો $m_R(\vec x, y) = y$ હોય અને $R(\vec{x}, y)$ લાગુ પડે}\\
y+1 & \text{અન્યથા.}
\end{cases}
\end{align*}
નોંધો કે $R(\vec x, z)$ લાગુ પડે એવી $z<y$ છે ત્યારે અને માત્ર
ત્યારે જ $m_R(\vec x,
y) \neq y$ છે.
\end{proof}

\begin{prob}
ધારો કે $R(\vec x, z)$ આદિમ પુનરાવર્તી છે. વિધેય
$m'_R(\vec{x}, y)$ને $\Char{R}$માંથી આદિમ પુનરાવર્તન વડે
વ્યાખ્યાયિત કરો. $R(\vec{x}, z)$ લાગુ પડે એવી કોઈ સંખ્યા $y$થી
નાનો હોય, તો તે સૌથી નાનો એવો~$z$ આપે અને અન્યથા $0$ આપે.
\end{prob}

\end{document}
